problem:  
Let \(V^n \subset S^{n+k}\) be a stationary integral \(n\)-varifold whose regular part is a smooth minimal submanifold, and assume that at every singular point \(p \in \mathrm{sing}\,V\) there exists a **unique tangent cone** \(C_p = \mathbb{R}_+ \times L_p\), where \(L_p \subset S^{n+k-1}\) is a smooth *stable* minimal \((n-1)\)-submanifold (the **link**).  
Denote by  
\[
\mathcal{L}(V) = \{ [L_p] : p \in \mathrm{sing}\,V \}
\]
the collection of isometry classes of these links, regarded as a subset of \(\mathrm{Var}_{n-1}(S^{n+k-1})\), the space of stationary integral \((n-1)\)-varifolds in the unit sphere equipped with the weak varifold topology.

**Question:**  
Assume all links \(L_p\) are stable and have uniformly bounded Morse index:
\[
\mathrm{ind}(L_p) \le N.
\]
Under these assumptions, is \(\mathcal{L}(V)\) necessarily **precompact in the smooth topology** (i.e. after appropriate ambient rotations, do the \(L_p\) subconverge in \(C^\infty\) to a minimal submanifold of \(S^{n+k-1}\))?  

Equivalently, can degeneration or bubbling of links occur, producing singular or non-smooth limits even under stability and bounded index?  
If such degenerations are possible, do they correspond precisely to local accumulation phenomena of singularities in \(V\)?  

In other words, does the combination of positive ambient curvature, stability, and finite index enforce *strong compactness* of tangent cone links, or can instability of geometric moduli still occur despite analytic control?

Why is it a "good" problem:  
This refined question corrects the dimensional setting and eliminates trivial compactness by moving from weak varifold topology to the **smooth topology under fixed index bounds**, a regime where standard varifold compactness no longer guarantees convergence. It asks whether geometric stability and bounded spectral complexity are sufficient to prevent degeneration of tangent cone links—a question at the frontier of **analytical and geometric compactness theory** for minimal varieties in spheres.

A positive answer would establish a powerful new rigidity link between **spectral control** (index bounds) and **geometric compactness**, implying that singularities can only exhibit finitely many geometric “link types.” A counterexample would reveal hidden analytic flexibility in minimal stability theory, illuminating the boundary between global spectral bounds and local geometric degeneration.

The problem combines tools from geometric measure theory (tangent cone analysis), differential geometry (minimal submanifolds in spheres), and analysis (bounded-index compactness), while remaining conceptually simple and geometric. Its resolution would produce a new form of **index–compactness principle**, a deep bridge between singular structure and spectral geometry—making it a “narrow entrance, profound depth” research problem.